\documentclass[11pt]{article}

\usepackage[a4paper,margin=1in]{geometry}
\usepackage[T1]{fontenc}
\usepackage[utf8]{inputenc}
\usepackage{lmodern}
\usepackage{amsmath,amssymb,amsthm}
\usepackage[hidelinks]{hyperref}
\usepackage{booktabs}
\usepackage{xcolor}

\title{Set-Period vs.\ Multiset-Period: Computational Results}
\author{Dirk Kunert}
\date{\today}

\begin{document}

\maketitle

\section*{Setup}

For each row $(o_n, o_d, a_n, a_d, \mathrm{period})$ of the file \\
\texttt{tests/new\_find\_patterns\_x\_max\_1000000\_51012\_lines.csv}, \\
the
program \texttt{src/c/set/add\_period\_set} computed the set-based period
length $\lambda_{\text{set}}$ on the same parameter quadruple and appended
it as a sixth column~\texttt{period\_set}. The run was performed in
\emph{global} mode with $X_{\max}=10^{6}$ and a per-row wall-clock cap of
$120$~seconds. The output is written to
\texttt{tests/new\_find\_patterns\_x\_max\_1000000\_51012\_lines.csv.tmp}.

\section*{Aggregate counts}

Of the $51{,}011$ data rows in the input, the resulting file breaks down
as follows.

\begin{center}
\begin{tabular}{lrl}
\toprule
Outcome & Rows & Meaning \\
\midrule
$\lambda_{\text{set}} = \mathrm{period}$  & $50{,}968$ & set period equals multiset period \\
$\lambda_{\text{set}} = -2$ & $29$ & \texttt{ARRAY\_SIZE\_EXCEEDED} \\
$\lambda_{\text{set}} = -3$ & $1$  & \texttt{DX\_LENGTH\_TO\_SMALL} \\
$\lambda_{\text{set}} = -4$ & $13$ & per-row $120$\,s timeout \\
$\lambda_{\text{set}} > 0$, $\lambda_{\text{set}} \neq \mathrm{period}$ & $1$ & genuine disagreement (see below) \\
\midrule
Total                                     & $51{,}012$ & \\
\bottomrule
\end{tabular}
\end{center}

In other words, on $99.92\%$ of the rows ($50{,}968 / 51{,}011$) the
set-based and multiset-based period lengths coincide; the remaining
$43$ rows split into $43$ computational aborts (out of memory, dx-window
too small, or timeout) and exactly one row in which the two periods are
both well-defined positive integers but disagree.

\paragraph{Interpretation under Theorem 6.1.}
The high agreement rate is structural, not coincidental. The input file
\texttt{new\_find\_patterns\_x\_max\_1000000\_51012\_lines.csv} is
concentrated in the regime $N < D$, in which Theorem~6.1 of the
companion paper
\texttt{rational\_cut\_and\_project\_gap\_periods.tex} forces
$\lambda_{\text{set}} = \lambda_{\text{multiset}}$.  The lone positive
disagreement therefore must lie in the complementary regime $N \ge D$,
where the theorem predicts $\lambda_{\text{set}} = 1$ — and indeed it
does (see the next two sections).

\section*{The single positive disagreement}

The only row with $\lambda_{\text{set}} > 0$ and
$\lambda_{\text{set}} \neq \mathrm{period}$ is

\begin{center}
\begin{tabular}{rrrrrr}
\toprule
$o_n$ & $o_d$ & $a_n$ & $a_d$ & multiset period & $\lambda_{\text{set}}$ \\
\midrule
$98{,}281$ & $139$ & $68$ & $149$ & $153{,}431$ & $1$ \\
\bottomrule
\end{tabular}
\end{center}

Here the multiset period is $153{,}431$ while the set-based period
equals $1$. The next section shows that $\lambda_{\text{set}} = 1$ is
the \emph{correct mathematical value} for this row, predicted by
Theorem~6.1 of the companion paper
\texttt{rational\_cut\_and\_project\_gap\_periods.tex}: the parameters
lie in the regime $N \ge D$, where the theorem forces
$\lambda_{\text{set}} = 1$.

\section*{Why $\lambda_{\text{set}} = 1$ is correct for this row}

The reported value is the mathematical answer predicted by
Theorem~6.1 of
\texttt{rational\_cut\_and\_project\_gap\_periods.tex}, and the
heuristic in \texttt{src/c/set/mathematics.c} arrives at it correctly.
Both perspectives are reviewed below.

\paragraph{Theorem 6.1 prediction.}
For this row $\alpha = 68$, $\beta = 149$, and $\omega = \gamma/\delta
= 98{,}281/139 \approx 707.06$. Hence
\[
D = \alpha^2 + \beta^2 = 4{,}624 + 22{,}201 = 26{,}825,
\]
\[
N = \lfloor\omega\alpha\rfloor + \lfloor\omega\beta\rfloor + 1
  = 48{,}079 + 105{,}351 + 1 = 153{,}431,
\]
so $N/D \approx 5.72 > 1$, i.e.\ $N \ge D$. Theorem~6.1 then yields
$\lambda_{\text{set}} = 1$, exactly the value reported.

\paragraph{Why every integer is a projected value.}
The condition $N \ge D$ is equivalent to: every residue class modulo
$D$ is hit by at least one accepted lattice point. Hence the
underlying set of projected values $\beta x + \alpha y$ is all of
$\mathbb{Z}$, and the gap sequence between consecutive distinct values
is identically~$1$.

\paragraph{Parameters and window.}
With $\gcd(\alpha,\beta)=1$ and $\gcd(\gamma,\delta)=1$ as above, the
width of the cut window in $y$ is
\[
\frac{\gamma(\alpha+\beta)}{\beta\,\delta}
= \frac{98{,}281 \cdot 217}{149 \cdot 139}
= \frac{21{,}326{,}977}{20{,}711} \approx 1029.74,
\]
so each $x$ contributes roughly $1030$ valid integer $y$'s — a very wide
strip. The actual run picked $x_{\max} = 1000$, producing $1{,}029{,}738$
raw projected values $\beta x + \alpha y$. This wide-strip regime is
exactly what makes $N \ge D$ here.

\paragraph{Sorted-difference structure.}
After sorting and dropping zero gaps (the set-collapse step in
\texttt{lambda} when \texttt{sort = true}), the sequence of distinct
sorted differences has $237{,}813$ entries:
\begin{itemize}
\item $232{,}589$ entries are equal to $1$ ($\approx 97.8\%$),
\item only $5{,}224$ entries are different from $1$ (and they take values
      in the small set \\ $\{2,3,4,7,10,13,16,29,42,55,68\}$),
\item the non-$1$ entries are concentrated entirely in the first
      $\sim\!6000$ and the last $\sim\!6000$ positions; the middle
      $95\%$ of the array — positions $5971$ through $231{,}821$ — is
      one uninterrupted run of $1$'s.
\end{itemize}
Because $\gcd(\alpha,\beta) = 1$ and the cut strip is wide, the
projected lattice $\beta x + \alpha y$ \emph{saturates} the integer line
in the bulk: consecutive distinct projected values differ by $1$ across
the entire interior. Only the two ends of the sorted dx-range — where
the strip ``runs out'' — produce gaps larger than $1$.

\paragraph{The trim-and-spiral heuristic locks onto the bulk.}
The period-detection logic in \\ \texttt{src/c/set/mathematics.c} starts by
trimming \texttt{(1 - FRACTION\_OF\_REMAINING\_ELEMENTS) / 40} of the
sequence from each end (here $594$ positions per side) and then spirals
both ends inward by $1$ each iteration, calling
\texttt{find\_period\_length} on the surviving window, until either a
period is detected or the window has shrunk by more than
$1 - \texttt{FRACTION\_OF\_REMAINING\_ELEMENTS} = 10\%$ of its initial
size. \texttt{find\_period\_length} returns $1$ exactly when every
element of the window equals every other element.

For this row the spiral reaches the all-$1$ bulk after $k \approx 5397$
inward steps: at $k = 5397$ the window is $[5991, 231{,}821]$, which lies
entirely inside the bulk run $[5971, 231{,}821]$ and is therefore
constant equal to $1$. The algorithm reports $\lambda_{\text{set}} = 1$
and stops.

\paragraph{Interpretation.}
The bulk of the sorted-and-collapsed difference sequence is the constant
$1$ because the projected lattice $\beta x + \alpha y$ \emph{saturates}
$\mathbb{Z}$: for every integer $\tilde p$ in a long central range, at
least one accepted lattice point projects to it. The boundary patches
at indices $[0,\,5970]$ and $[231{,}822,\,237{,}812]$ are finite-window
artefacts of computing on a bounded $x$-range; in the bi-infinite
sequence they vanish and the sequence is the constant~$1$ everywhere.
This saturation is exactly the geometric content of the case
$N \ge D$ in Theorem~6.1, so the heuristic's
report $\lambda_{\text{set}} = 1$ is mathematically correct, not an
artefact.

\section*{Run summary}

\begin{itemize}
\item Resume from a partial run: $40{,}201$ rows already present in the
      output were skipped; $10{,}811$ rows were freshly processed.
\item No hard failures during the resume run; $11$ of the $13$ timeouts
      shown above were produced in this run, the remaining $2$ were
      already present in the partial output.
\item Errors of types $-2$ and $-3$ accumulated $30$ rows in total;
      these are computational limitations of the current
      \texttt{lambda} implementation at $X_{\max}=10^{6}$ and do not
      indicate a true mismatch of the two period notions.
\end{itemize}

\section*{Caveats}

\begin{itemize}
\item Negative \texttt{period\_set} values mean the set-based period was
      \emph{not} computed, not that it differs from the multiset period.
      Excluding the $43$ aborted rows, the multiset period and the
      set-based period agree on every single one of the remaining
      $50{,}968$ rows except the row exhibited above.
\item The timeout cap of $120$\,s is heuristic; lifting it (or doubling
      \texttt{MAX\_PERIOD\_ARRAY\_SIZE}) might convert some of the $43$
      aborts into successful computations.
\end{itemize}

\section*{Re-check of the $43$ aborted rows}

The $43$ rows that produced negative \texttt{period\_set} values in the
first pass were extracted into
\texttt{tests/errors\_set\_negative\_lambda.csv} and re-run with three
limits raised:
\texttt{MAX\_PERIOD\_ARRAY\_SIZE} from $10^{8}$ to $4\!\cdot\!10^{9}$
(allocating $\approx 32$~GB for the \texttt{dx} buffer),
\texttt{X\_MAX} from $10^{6}$ to $10^{9}$, and the per-row wall-clock cap
to $180$~s. Output: \texttt{tests/errors\_set\_negative\_lambda\_recheck.csv}.

\begin{center}
\begin{tabular}{lrl}
\toprule
Outcome & Rows & Meaning \\
\midrule
$\lambda_{\text{set}} = \mathrm{period}$ & $26$ & newly resolved, set period equals multiset period \\
$\lambda_{\text{set}} = -4$              & $17$ & still timed out at $180$~s per retry slice \\
$\lambda_{\text{set}} > 0$, $\lambda_{\text{set}} \neq \mathrm{period}$ & $0$ & no new disagreement \\
\midrule
Total                                    & $43$ & \\
\bottomrule
\end{tabular}
\end{center}

Crucially, every one of the $26$ rows newly resolved by the larger
buffer agrees with the multiset period: the set-based and multiset-based
period coincide on all of them. The remaining $17$ rows are not
disagreements; they are computations that the current implementation
cannot complete in a $180$-second budget per retry, because the retry
loop in \texttt{add\_period\_set.c} doubles \texttt{target\_points} each
time it receives \texttt{DX\_LENGTH\_TO\_SMALL} or \texttt{NO\_PERIOD}
and starts a fresh alarm window.

\paragraph{Theorem 6.1 resolution of the 17 timeout rows.}
Direct computation of $N = \lfloor\omega\alpha\rfloor +
\lfloor\omega\beta\rfloor + 1$ and $D = \alpha^2 + \beta^2$ on the
parameter quadruples of the $17$ remaining timeout rows shows that
\emph{every one of them satisfies $N < D$}, with $N/D$ ranging from
about $10^{-4}$ to $0.4$. Theorem~6.1 then forces
$\lambda_{\text{set}} = N = \lambda_{\text{multiset}}$ for all $17$
rows, where the multiset period is already known from the input file.
No actual set-period computation is needed: the answer is mathematically
determined by the multiset period, and the $17$ timeouts represent
redundant work that the C heuristic could skip entirely on input rows
with $N < D$. Redesigning \texttt{add\_period\_set.c} to short-circuit
on $N < D$ (read the multiset period directly) is a cheap implementation
improvement but is out of scope for this report.

\section*{Updated picture for the full 51{,}011-row dataset}

Combining the original pass and the re-check:

\begin{center}
\begin{tabular}{lr}
\toprule
Outcome & Rows \\
\midrule
$\lambda_{\text{set}} = \mathrm{period}$ (confirmed agreement) & $50{,}994$ \\
$\lambda_{\text{set}} > 0$, $\lambda_{\text{set}} \neq \mathrm{period}$ (positive disagreement) & $1$ \\
$\lambda_{\text{set}} = -4$ (still inconclusive at $180$~s/retry) & $17$ \\
$\lambda_{\text{set}} \in \{-1,-2,-3\}$ (still inconclusive otherwise) & $0$ \\
\midrule
Total                                                            & $51{,}012$ \\
\bottomrule
\end{tabular}
\end{center}

This raises the confirmed-agreement rate to
$50{,}994 / 51{,}011 \approx 99.97\%$, with the lone positive
disagreement
$(o_n,o_d,a_n,a_d) = (98{,}281,\,139,\,68,\,149)$,
$\mathrm{period} = 153{,}431$, $\lambda_{\text{set}} = 1$
unchanged from before. None of the $26$ newly resolved rows produced an
additional positive disagreement.

Under Theorem~6.1 the picture is structurally complete. Direct
computation of $N$ and $D$ on every row shows that exactly one row
of the $51{,}012$ inputs lies in the regime $N \ge D$ (the lone
disagreement above), and all $51{,}011$ remaining rows satisfy $N < D$.
The theorem then dictates: the $50{,}994$ confirmed agreements are
forced by $N < D$ (set and multiset coincide), the lone disagreement
gives $\lambda_{\text{set}} = 1$ as predicted by the $N \ge D$ branch,
and the $17$ remaining timeouts also lie in $N < D$, so their
set-period coincides with the multiset period already in the input file
(see the post-mortem below). No further input row can produce a fresh
positive disagreement that contradicts the theorem.

\end{document}
